%% ===========================================================================
\FloatBarrier
\section{Formulation}
\label{sec:formulation}
%% ===========================================================================

The model is a two-dimensional electromechanical continuum with three
components: the adopted strain-gradient energy is parametrized by the second
moment of a prescribed rectangular window, the electric enthalpy includes the
piezoelectric and the flexoelectric couplings together, and the reduced mechanical response is assigned a phenomenological
fractional Kelvin--Voigt stiffness-memory law. The window geometry, kinematics, energy, constitutive equations,
two-dimensional reduction and fractional memory law are introduced
before deriving the governing equations. Closed-form coefficients for the prescribed window shapes are given in
\ref{app:Mcoeff}.

\FloatBarrier
\subsection{Nonlocal averaging window and the second-moment tensor}
\label{sec:window}

Repeated indices are summed. Let the strain gradient $\eta_{ijk}$ at an
in-plane position $\bm{x}$ be replaced by its average
over a window $\mathcal{W}$ of area $A_{\mathcal W}$ centered at that point,
\begin{equation}
  \bar\eta_{ijk}(\bm{x}) = \frac{1}{A_{\mathcal W}}
     \int_{\mathcal W} \eta_{ijk}(\bm{x}+\bm{s})\,\mathrm{d}\bm{s}.
  \label{eq:average}
\end{equation}
Expanding the integrand to second order in the window coordinate $\bm{s}$,
\begin{equation}
  \eta_{ijk}(\bm{x}+\bm{s}) \simeq \eta_{ijk}(\bm{x})
     + s_m\,\eta_{ijk,m}(\bm{x}) + \tfrac12 s_m s_n\,\eta_{ijk,mn}(\bm{x}),
  \label{eq:expand}
\end{equation}
and integrating term by term, the linear term vanishes because the window is
centered, so
\begin{equation}
  \bar\eta_{ijk} \simeq \eta_{ijk} + \tfrac12 M_{mn}\,\eta_{ijk,mn},
  \qquad
  M_{mn} = \frac{1}{A_{\mathcal W}}\int_{\mathcal W} s_m s_n \,\mathrm{d}\bm{s},
  \label{eq:Mdef}
\end{equation}
where the retained geometric information enters the truncated model through
the symmetric second-order tensor $M$, which has the dimensions of a squared
length and plays the role that a scalar characteristic length plays in
isotropic strain-gradient elasticity. For the windows compared in this work, the diagonal components follow from elementary integration.
The rectangle and ellipse are averaged with the two-dimensional area measure
in \eqref{eq:average}; the cuboid and ellipsoid are three-dimensional
volume averages, whose in-plane moments are imported as parameters into
the reduced two-dimensional model:
\begin{equation}
  M^{\rm cuboid}_{ii} = \frac{a_i^2}{3},\qquad
  M^{\rm ellipse}_{ii} = \frac{a_i^2}{4},\qquad
  M^{\rm ellipsoid}_{ii} = \frac{a_i^2}{5},
  \label{eq:Mshapes}
\end{equation}
where $a_i$ denote the half-widths of rectangles/cuboids and the semi-axes
of ellipses/ellipsoids. The respective integration measures are
$\mathrm{d}^2 s$ and $\mathrm{d}^3 s$; the resulting moments are
$a_i^2/(d+2)$ for a uniform $d$-dimensional ellipse/ellipsoid.
A window rotated by $\theta$ in the plane is obtained from the axis-aligned
tensor by $M\mapsto RMR^{\mathsf T}$, that is
\begin{equation}
  M_{11}(\theta)=M_{11}\cos^{2}\theta+M_{22}\sin^{2}\theta,\qquad
  M_{12}(\theta)=\tfrac12\left(M_{11}-M_{22}\right)\sin2\theta,
  \label{eq:Mrot2}
\end{equation}
so the signed off-diagonal component varies with $\sin2\theta$ and the diagonal ones
carry the orientation dependence.

A consequence of \eqref{eq:Mshapes}, examined numerically below, is that at fixed $M$ the
\emph{shape} of the window carries no additional information within the
adopted second-moment approximation, because two windows with the same second
moment produce identical solutions of this truncated model. This does not imply
equivalence of the original integral kernels; in the present approximation window anisotropy, specified by aspect ratio and orientation, enters
the adopted material model through $M$, since it
changes $M$ and with it the ratio between the directional gradient penalties.

\FloatBarrier
\subsection{Kinematics}
\label{sec:kinematics}

Let $\Omega\subset\mathbb{R}^2$ be the mid-plane of a strip of thickness $h$,
with length $L$, coordinates $x_1$ along the member and $x_2$ through its thickness. The
mechanical state is the pair $(u_1,u_2)$, where $u_1$ is the in-plane
displacement and $u_2\equiv w$ is the transverse deflection, the model retains
the full second gradient of both components, and the strain and the strain
gradient are
\begin{equation}
  \varepsilon_{ij} = \tfrac12\left(u_{i,j}+u_{j,i}\right),\qquad
  \eta_{ijk} = \partial_k \varepsilon_{ij} = \tfrac12\left(
  u_{i,jk}+u_{j,ik}\right),
  \label{eq:strain}
\end{equation}
so that the mechanical degrees of freedom at a node are the eight quantities
\begin{equation}
  \left\{u_1,\; \partial_1u_1,\; \partial_2u_1,\; \partial_{12}u_1,\;
         u_2,\; \partial_1u_2,\; \partial_2u_2,\; \partial_{12}u_2\right\},
  \label{eq:dofs}
\end{equation}
that is, the value, the two first derivatives and the cross derivative of each
component -- eight mechanical degrees of freedom per node and not a
three-component continuum. These are the degrees of freedom of the
Bogner--Fox--Schmit rectangle~\cite{BognerFoxSchmit1966}: a bicubic Hermite
tensor-product element, $C^1$ across element edges and $C^0$ only for the
cross derivative, in which the displacement of one component is interpolated
from its nodal data as
\begin{equation}
  u(\xi,\eta) = \sum_{a=1}^{4}\Big[N^{00}_{a}u_a + N^{10}_{a}u_{,x}\big|_a
    + N^{01}_{a}u_{,y}\big|_a + N^{11}_{a}u_{,xy}\big|_a\Big],
  \label{eq:interp}
\end{equation}
with $N^{pq}_a$ the products of the one-dimensional Hermite cubics. This interpolation represents the second displacement gradients in two
dimensions and was used in the one-scale theory~\cite{fem1}. The present
formulation incorporates flexoelectric coupling and fractional memory.

\FloatBarrier
\subsection{Electric enthalpy and the piezoelectric--flexoelectric
constitutive relations}
\label{sec:energy}

The electromechanical enthalpy of the strip is assembled from five terms. The
elastic energy is
\begin{equation}
  \Pi_{\rm el} = \tfrac12\int_\Omega \varepsilon_{ij} C_{ijkl}
        \varepsilon_{kl}\,\mathrm{d}A ,
  \label{eq:Eel}
\end{equation}
the adopted strain-gradient energy is parameterized by the window second
moment defined in \eqref{eq:Mdef}. This is a separately postulated positive
local constitutive closure. In particular, expanding an ordinary quadratic
energy under a normalized spatial-average window would give a variance
contribution $\tfrac12 M_{mn}\varepsilon_{ij,m}C_{ijkl}\varepsilon_{kl,n}$,
not the coefficient chosen below. Thus the $1/4$ is a model normalization,
not a consequence of that elementary variance expansion. In the adopted model,
\begin{equation}
  \Pi_{\rm gr} = \tfrac14\int_\Omega M_{mn}\,
        \varepsilon_{ij,m}\,C_{ijkl}\,\varepsilon_{kl,n}\,\mathrm{d}A ,
  \label{eq:Egrad}
\end{equation}
the piezoelectric and dielectric energies are
\begin{equation}
  \Pi_{\rm pz} = -\int_\Omega e_{kij} E_k \varepsilon_{ij}\,\mathrm{d}A
  \qquad\text{and}\qquad
  \Pi_{\rm dl} = -\tfrac12\int_\Omega \kappa_{ij} E_i E_j\,\mathrm{d}A ,
  \label{eq:Epiezo}
\end{equation}
and the flexoelectric term is the polarization-gradient coupling of Mindlin's
Form~II~\cite{Mindlin1965,Toupin1962,Mindlin1968},
\begin{equation}
  \Pi_{\rm fx} = -\int_\Omega \mu_{kijl} E_k \eta_{ijl}\,\mathrm{d}A ,
  \label{eq:Eflexo}
\end{equation}
where $E_k=-\phi_{,k}$ is the electric field, $\phi$ the electric potential,
$C_{ijkl}$ the elastic stiffness, $e_{kij}$ the piezoelectric tensor,
$\kappa_{ij}$ the permittivity, $\mu_{kijl}$ the flexoelectric tensor,
$W_{\rm ext}$ the work of the external loads. The total enthalpy is their sum,
\begin{equation}
  \Pi = \Pi_{\rm el} + \Pi_{\rm gr} + \Pi_{\rm pz} + \Pi_{\rm dl}
        + \Pi_{\rm fx} - W_{\rm ext},
  \label{eq:energy}
\end{equation}
where $W_{\rm ext}$ denotes the sum of applied body work and
prescribed mechanical or electrical boundary work. These contributions
are separated in the weak form; a prescribed surface load must not be
counted once in $W_{\rm ext}$ and again as an additional right-hand side
term. In this two-dimensional formulation all integrals represent
energy per unit out-of-plane width. Here the isotropic length
$\ell^2\delta_{mn}$ of a scalar theory is
replaced by $M_{mn}$ in \eqref{eq:Egrad}. The flexoelectric tensor is restricted to a one-parameter isotropic
specialization and is represented by the single coefficient
\begin{equation}
  \mu_{kijl} = \mu\,\delta_{ij}\delta_{kl},
  \qquad \mu = 10^{-6}\ \si{\coulomb\per\meter},
  \label{eq:muexplicit}
\end{equation}
a charge per unit length, prescribed for a numerical sensitivity study.
This coefficient is not inferred from the dielectric/piezoelectric data of
Zgonik et al., and it is not a specimen-specific measured value.
The measured transverse response of unpoled ceramic BaTiO$_3$
($\mu_{12}\simeq\SI{10}{\micro\coulomb\per\meter}$)~\cite{MaCross2006}
and the longitudinal shock-wave response of unpolarized bulk BaTiO$_3$
($\mu_{11}\simeq\SI{17.33}{\micro\coulomb\per\meter}$)~\cite{HuDengLiangShen2017}
are useful experimental context, but cannot be substituted for this parameter.
Indeed, $\mu_{kijl}=\mu\delta_{ij}\delta_{kl}$ couples the electric
field to the gradient of the in-plane strain trace, rather than isolating a
single experimentally measured component. A component-matched tensor,
phase, poling state, specimen and measurement regime would be needed for
quantitative calibration. The flexoelectric part of the electromechanical
coupling matrix is linear in $\mu$; a flexoelectric-only coupling factor has
a leading quadratic dependence on $\mu$ for weak coupling, whereas mixed
piezoelectric--flexoelectric configurations also contain cross terms.

Let $\mathcal H(\varepsilon,\nabla\varepsilon,\bm E)$ be the local
electromechanical enthalpy density, namely the sum of the five integrands
in \eqref{eq:Eel}--\eqref{eq:Eflexo} before domain integration or
subtraction of external work. The constitutive relations are local partial
derivatives of $\mathcal H$. The stress is
\begin{equation}
  \sigma_{ij} = \frac{\partial \mathcal H}{\partial \varepsilon_{ij}}
              = C_{ijkl}\,\varepsilon_{kl} - e_{kij} E_k ,
  \label{eq:sig}
\end{equation}
the higher-order stress is
\begin{equation}
  \varsigma_{ijm} = \frac{\partial \mathcal H}{\partial \varepsilon_{ij,m}}
              = \tfrac12 M_{mn}\,C_{ijkl}\,\varepsilon_{kl,n}
                - \mu_{kijm}\,E_k ,
  \label{eq:tau}
\end{equation}
and the electric displacement is
\begin{equation}
  D_k = -\frac{\partial \mathcal H}{\partial E_k}
              = e_{kij}\,\varepsilon_{ij} + \kappa_{kl}\,E_l
                + \mu_{kijl}\,\eta_{ijl},
  \label{eq:D}
\end{equation}
Thus the electric displacement depends on strain, electric field and
strain gradient, while the higher-order stress includes both gradient-elastic
and flexoelectric contributions. With the single-constant tensor \eqref{eq:muexplicit} the last two
relations become
\begin{equation}
  \varsigma_{ijm} = \tfrac12 M_{mn}\,C_{ijkl}\,\varepsilon_{kl,n}
                    - \mu E_m\delta_{ij},
  \qquad
  D_k = e_{kij}\varepsilon_{ij} + \kappa_{kl}E_l
        + \mu\,\partial_k\left(\varepsilon_{xx}+\varepsilon_{yy}\right),
  \label{eq:Dexplicit}
\end{equation}
which is used in the finite-element formulation. An ellipsoidal window carries
$M_{mn}=L_{mn}/5$ in the notation of the one-scale theory with a tensorial
characteristic length, so $\tfrac14 M_{mn}=\tfrac1{20}L_{mn}$ and the gradient
energy of~\cite{fem1} is recovered exactly.

The three relations are derivatives of the local density $\mathcal H$
associated with \eqref{eq:energy}, evaluated with $E_k=-\phi_{,k}$, and their signs are those of the implementation: the
dielectric energy enters $\Pi$ with a minus sign so that $D$ carries
$+\kappa_{kl}E_l$, the assembled dielectric block is
$K_{\phi\phi}=+\int\nabla N_\phi^{\mathsf T}\kappa\nabla N_\phi$, the
piezoelectric coupling is $B=\int B^{\mathsf T}e^{\mathsf T}\nabla N_\phi$, and
the flexoelectric coupling $+\int\mu\,\nabla N_\phi^{\mathsf T}
\nabla(\varepsilon_{xx}+\varepsilon_{yy})$ is carried by the same gradient of
the potential. The transpose structure of the electromechanical blocks follows from
the common enthalpy.

The averaging identity \eqref{eq:Mdef} determines the window second moment,
but does not by itself derive the positive local gradient-energy term
\eqref{eq:Egrad} or its prefactor: averaging the local elastic energy also
produces terms involving second derivatives of strain and boundary
contributions. Equation~\eqref{eq:Egrad} is therefore the specified
second-moment constitutive closure, with the factor $1/4$ taken as part of
the model definition, not an exact identity for an arbitrary integral
nonlocal material. Published mixed formulations of flexoelectricity
introduce higher-order moduli as constitutive data and do not establish
this specific window-to-modulus normalization~\cite{MaoPurohitAravas2016,SerraoKozinov2024}.
Its applicability requires a separate comparison with
the underlying nonlocal response or experiment (see Sec.~\ref{sec:limit_window}); $a_i/L$ alone cannot establish that validity.

Equation~\eqref{eq:energy} has three implications for the discretization.
First, the
gradient energy depends on the second derivatives of the displacement, so the element
must be $C^1$ in the sense of \eqref{eq:dofs}; a $C^0$ element would make the
fourth-order term non-conforming. Second, the piezoelectric and flexoelectric
couplings are included in the same enthalpy but contribute differently under the
open- and short-circuit electrode configurations.
Third, because the strip is treated as a two-dimensional continuum, the
plane-stress or plane-strain reduction of $C$, $e$ and $\kappa$ is part of the
model rather than a detail of the implementation.

\FloatBarrier
\subsection{Two-dimensional reduction of the material tensors}
\label{sec:reduction}

The strip is free out of plane, so the stress component normal to the plane
vanishes and the stiffness that multiplies the in-plane strain is the reduced
tensor
\begin{equation}
  C^{p}_{ij} = C_{ij} - \frac{C_{i3}C_{j3}}{C_{33}},
  \label{eq:Cp}
\end{equation}
with the corresponding reductions of the piezoelectric and dielectric tensors,
\begin{equation}
  e^{p}_{ai} = e_{ai} - \frac{e_{a3}C_{3i}}{C_{33}},
  \qquad
  \kappa^{p}_{ab} = \kappa_{ab} + \frac{e_{a3}e_{b3}}{C_{33}},
  \label{eq:ep}
\end{equation}
where $a,b$ label in-plane electric-field components and $i$ labels the
in-plane normal-strain components in Voigt notation. The shear coupling is
unchanged in this reduced crystal orientation. The plane-strain alternative
keeps the unreduced in-plane blocks and imposes
$\varepsilon_{zz}=0$. These algebraic formulae refer to the \emph{local}
piezoelectric constitutive blocks; a consistent three-dimensional reduction
of the full strain-gradient and flexoelectric model would additionally require
the associated higher-gradient boundary conditions and is not claimed here.

The two reductions are not interchangeable: for the length-poled strip, the plane-stress entry is
$C^{p}_{xx}=9.5500\times10^{10}$\,Pa against $C_{xx}=1.51\times10^{11}$\,Pa in
plane strain, and the extensional coupling factor of the two is $0.5727$ and
$0.5132$ respectively. The plane-stress reduction is used everywhere below
because the strip is free out of plane. The difference between the two
reductions is quantified in Sec.~\ref{sec:verification}.

\FloatBarrier
\subsection{Fractional Kelvin--Voigt memory}
\label{sec:fracmem}

The reduced mechanical model is assigned a phenomenological fractional
Kelvin--Voigt memory on a chosen assembled stiffness matrix. The local
constitutive modulus and the assembled FEM stiffness matrix are distinct
objects with different dimensions and units. The displayed electrical
displacement law \eqref{eq:Dexplicit} is instantaneous: it does not include
a fractional derivative of the electric field or the permittivity. For a
sufficiently regular function, the Caputo derivative is
\begin{equation}
  {}^{C}\!D_t^{\alpha}f(t) = \frac{1}{\Gamma(1-\alpha)}
     \int_{0}^{t}\frac{\dot f(t')}{(t-t')^{\alpha}}\,\mathrm{d}t',
  \qquad 0<\alpha<1,
  \label{eq:caputo}
\end{equation}
and the spring-pot of Caputo~\cite{Caputo1967}, whose use as a constitutive
element is due to Bagley and Torvik~\cite{BagleyTorvik1983} and
Koeller~\cite{Koeller1984}, is
\begin{equation}
  \sigma(t) = C\,\varepsilon(t) + C\,\tau^{\alpha}\,
              {}^{C}\!D_t^{\alpha}\varepsilon(t),
  \label{eq:kv}
\end{equation}
with $\tau$ a relaxation time and $C$ a scalar local constitutive
modulus. This scalar relation illustrates the fractional element; the
numerical simulations instead prescribe $\tau^\alpha\mathbf K_d
{}^{C}\!D_t^{\alpha}\mathbf u$ in the assembled generalized-coordinate
equation, with $\mathbf K_d$ explicitly selected from the
condensed or mechanical stiffness matrices. The constitutive law here is a fractional stiffness-memory model, not
a derivation of the Cole--Cole dielectric permittivity
\cite{ColeCole1941,HavriliakNegami1967}. Using the $\exp(+\mathrm{i}\omega t)$ harmonic convention, its
harmonic steady-state modulus is
\begin{equation}
  C^{*}(\omega) = C\left[1+(\mathrm{i}\omega\tau)^{\alpha}\right],
  \label{eq:kvfreq}
\end{equation}
under the principal branch for the complex power. The alternate
$\exp(-\mathrm{i}\omega t)$ convention replaces $\mathrm{i}\omega$
by $-\mathrm{i}\omega$. The transient program uses the Laplace variable
$s$ and retains the initial term
$\mathcal L\{{}^{C}D_t^\alpha q\}=s^\alpha Q-s^{\alpha-1}q(0)$.
The constitutive law reduces to the classical Kelvin--Voigt element for $\alpha=1$,
for which the Caputo operator is defined separately as the ordinary
first derivative.

In the condensed mechanical problem the postulated memory operator
$\mathbf K_d$ can be chosen as either $K_{\rm eff}$, the fully
condensed stiffness incorporating the electrical boundary condition, or
$K_{uu}$, the stiffness retained for the mechanical-only memory variant.
Define $\mathbf K_{\rm mem}=\tau^\alpha\mathbf K_d$, consistently with the time integrator. These are distinct phenomenological placements of memory, not a consequence
of identifying the local modulus $C$ with a global matrix; their predicted
damping is compared numerically below. Discretizing \eqref{eq:caputo} with the L1 scheme at uniform steps
$\Delta t$, with $n$ the step index and $t_n=n\Delta t$, gives
\begin{equation}
  {}^{C}\!D_t^{\alpha}\varepsilon(t_n) \approx
  \frac{1}{\Gamma(2-\alpha)}\sum_{k=1}^{n} b_{n-k}\,
  \frac{\varepsilon^{k}-\varepsilon^{k-1}}{\Delta t^{\alpha}},
  \label{eq:l1}
\end{equation}
with the weights
\begin{equation}
  b_j=(j+1)^{1-\alpha}-j^{1-\alpha},\qquad b_0=1,
  \label{eq:bdef}
\end{equation}
evaluated before time integration and reused at each step. The weights decay algebraically, retaining contributions from earlier
time steps, and at $\alpha=1$ they collapse to $b_0=1$, $b_j=0$ for $j\ge1$, so
that the scheme becomes a backward difference and the memory reduces to the
classical Kelvin--Voigt law up to its first-order truncation error. The
convolution is applied to the reduced modal coordinates, so the history that
has to be stored is one scalar sequence per retained mode and the cost per step
does not grow with the number of elements; which of the two matrices enters
\eqref{eq:kv} decides whether the piezoelectric contribution is inside the
remembered operator or outside it. For $0<\alpha<1$ and sufficiently smooth functions, the standard L1
formula has local truncation error $O(\Delta t^{2-\alpha})$ under its
usual regularity assumptions; weak initial-time regularity and the coupled
Newmark--L1 solver can give a lower observed order. At $\alpha=1$ the
limiting backward difference is first order. The selected step size and
observed convergence for the coupled solver are reported with the numerical
benchmarks.

\FloatBarrier
\subsection{Governing equations}
\label{sec:governing}

Hamilton's principle applied to the kinetic energy and to the
electromechanical enthalpy \eqref{eq:energy}, with Gauss' law, gives
the dynamic strong form on $\Omega$,
\begin{align}
  \sigma_{ij,j} - \varsigma_{ijm,jm} - \rho\,\ddot u_i &= 0,
  \label{eq:gov1}\\
  D_{k,k} &= 0,
  \label{eq:gov2}
\end{align}
where $\rho$ is the mass density. These equations hold together with the
classical traction and the higher-order double-traction
conditions on the boundary. With
\eqref{eq:Dexplicit} substituted, the two equations read
\begin{equation}
  \left(C_{ijkl}\varepsilon_{kl}-e_{kij}E_k\right)_{,j}
    - \left[\tfrac12 M_{mn}C_{ijkl}\varepsilon_{kl,n}\right]_{,jm}
    + \mu\left(\nabla\cdot\bm{E}\right)_{,i}
    - \rho\ddot u_i = 0,\qquad
  \left(\kappa_{kl}E_l\right)_{,k} =
    -\left(e_{kij}\varepsilon_{ij}\right)_{,k}
    -\mu\,\partial_k\partial_k
    \left(\varepsilon_{xx}+\varepsilon_{yy}\right),
  \label{eq:govexp}
\end{equation}
so that the strong form retains the piezoelectric terms in both
fields, the mass term has the sign required by the dynamic weak form,
and \eqref{eq:gov1} carries the flexoelectric term
$\mu(\nabla\cdot\bm E)_{,i}$ and \eqref{eq:gov2} carries
$\mu\partial_k\partial_k(\varepsilon_{xx}+\varepsilon_{yy})$; the two are the
transpose pair that the weak form assembles. The mechanical boundary condition
of the gradient theory is the pair $(t_i,\,m_i)$ -- ordinary traction [Pa] and double traction [N/m] per unit out-of-plane width -- so a clamped edge constrains both the displacement and its
normal derivative, exactly the quantities that appear in \eqref{eq:dofs}. In
the time domain the memory of \eqref{eq:kv} turns \eqref{eq:gov1} into a
Volterra-type equation whose kernel is the fractional derivative of
\eqref{eq:l1}; the discretized weak form is written directly for the
time-discrete problem, which is what the implementation solves.

The following parameters are varied separately in the numerical study: the memory of
\eqref{eq:kv} acts on the \emph{damping} at fixed eigenfrequency, the window of
\eqref{eq:Mdef} acts on the \emph{stiffness} at fixed memory, and the
flexoelectric tensor $\mu$ acts on the \emph{coupling} at fixed geometry. Unless stated otherwise, parameter comparisons hold the remaining
inputs fixed.

The electrode configurations are illustrated in Fig.~\ref{fig:geometry}: the bending family is poled across the thickness with
electrodes on the long faces, and the extensional family is poled along the
length with electrodes on the end faces. They are two electrical
configurations of one strip, and the coupling factors obtained from them are
therefore not additive.

The model parameters are summarized in Table~\ref{tab:params}. The
elastic, piezoelectric and dielectric constants follow
Table II of Zgonik and co-workers for room-temperature tetragonal
single-crystal $\mathrm{BaTiO_3}$~\cite{Zgonik1994}. That source also gives
$C^E_{1212}=\SI{134}{\giga\pascal}$ for basal-plane shear; it does not
enter the present in-plane shear calculation because each simulated plane
contains the polar axis and one basal axis, for which
$C^E_{1313}=\SI{61}{\giga\pascal}$ is used. The numerical values are
nominal fitted constants with experimental uncertainty, not universal
crystal-independent constants.
The prescribed mass density $\rho=6020$\,\si{\kilogram\per\cubic\meter}
is a rounded representative crystal density: single-crystal diffraction
at $298$~K gives a calculated crystallographic density of
$6019$\,\si{\kilogram\per\cubic\meter}~\cite{Nakatani2016}.
The difference is about $0.017\,\%$; this agreement supports the chosen
nominal inertial parameter, but does not establish a measured density for
the hypothetical strip or its microstructure. Mode frequencies and time
scales remain conditional on the density choice. The additional model and numerical inputs are prescribed rather than
measured, including the flexoelectric coefficient, taken as
$\mu=10^{-6}$\,\si{\coulomb\per\meter}, which is an assumed sensitivity value, not a
measurement of the specific specimen considered here. Its effect is reported for this assumed value; only the
isolated, weak-coupling flexoelectric contribution has a simple $\mu^2$
leading dependence. The window half-widths, constrained to the assumed range
$a/L\le0.2$ for this second-moment model;
the fractional order and relaxation time, which are free parameters of the
memory law and are varied over the ranges stated; the specimen geometry;
and the finite-element mesh, for which refinement comparisons accompany
the numerical benchmarks.

%% The parameter table: quantity, symbol, value, unit.  The source of every
%% value is stated in the sentence above the table, not in a column.
\begin{table}[!t]
\centering
\small
\caption{Parameters of the model. Elastic, piezoelectric and dielectric parameters are attributed to
the cited material reference; density is a rounded representative
crystallographic value~\cite{Nakatani2016}, not measured for this specimen;
flexoelectric coefficient, window, memory-law parameters, geometry and mesh
are prescribed inputs, not independently calibrated material measurements.}
\label{tab:params}
\begin{tabular}{@{}llll@{}}
\toprule
quantity & symbol & value & unit \\
\midrule
elastic constants & $C_{11}$, $C_{12}$, $C_{13}$, $C_{33}$, $C_{44}$ &
$222$, $108$, $111$, $151$, $61$ & GPa \\
piezoelectric constants & $e_{15}$, $e_{31}$, $e_{33}$ &
$34.2$, $-0.7$, $6.7$ & \si{\coulomb\per\square\meter} \\
relative permittivities & $\kappa_{11}/\kappa_0$, $\kappa_{33}/\kappa_0$
(polar axis) &
$2200$, $56$ & -- \\
mass density & $\rho$ & $6020$ & \si{\kilo\gram\per\cubic\meter} \\
\midrule
flexoelectric coefficient & $\mu$ & $10^{-6}$ & \si{\coulomb\per\meter} \\
window half-widths & $a_1$, $a_2$ & $3$, $1$ & \si{\micro\meter} \\
fractional order (studied range) & $\alpha$ & $0.2$--$1$ & -- \\
relaxation time (studied range) & $\tau$ & $0.3$--$32$ & \si{\nano\second} \\
strip geometry & $L$, $h$ & $200$, $10$ & \si{\micro\meter} \\
mesh & -- & $48\times4$ BFS elements & -- \\
\bottomrule
\end{tabular}
\end{table}


Limiting cases follow by suppressing individual constitutive terms.

\begin{figure}[!t]
\centering
\includegraphics[width=\textwidth]{F1_geometry.pdf}
\caption{The two electrode configurations and the averaging window.
(a) Bending family, poled through the thickness: electrodes on the large
faces, the lower earthed and the upper floating at zero net charge.
(b) Extensional family, poled along the length: electrodes on the short ends.
The window second moment fixes the gradient energy. The strip length is $L=\SI{200}{\micro\meter}$ and its thickness is $h$; the window half-widths are $a_1\times a_2$, with $M_{ii}=a_i^2/3$ for a cuboid and $a_i^2/4$ for an ellipse. The schematic is not to scale.}
\label{fig:geometry}
\end{figure}

\FloatBarrier
